% Fragmento integrable. Preambulo anfitrion: amsmath, amssymb, longtable, hyperref.
\section{Profil dodécaphasique et retour quadratique}
\label{hmtfg:fragmento:perfil-dodecafasico-y-retorno-cuadratico}

\subsection{Généalogie et objet de la démonstration}
\label{hmtfg:desarrollo:1}

L’objet initial est la lecture critique de la roue dodécaphasique orientée déjà construite dans le corpus. APP évalue la somme et le produit sur les deux feuilles du réseau de chiffres, en conservant quotient et reste ; TRIT détermine le régime et l’orientation ; le transport TPK sélectionne, transporte et actualise l’état avec ses registres de retenue, de frontière et de mémoire. La connexion nonadique prolonge cet état, avec retour de phase et avancement de mémoire. Ses lecteurs agissent sur la même structure discrète du continu.

Le lecteur focal de la source est
\[
\mathcal P_{\mathrm{crit}}:U_{12}^{\mathrm{sign}}\longrightarrow
\mathbb R/2\pi\mathbb Z,\qquad
m\longmapsto\phi_m=\frac{(3m-1)\pi}{18},\quad 1\le m\le12.
\]
Pour évaluer le moment angulaire, on conserve le représentant orienté de cette carte : le remplacer par un autre représentant modulo \(2\pi\) modifierait le moment. La trace uniforme fixe \(w_m=1/12\). La chaîne focale est
\[\begin{aligned}
\text{APP–TRIT–TPK}&\longrightarrow
\text{structure discrète du continu}\\ &\longrightarrow U_{12}^{\mathrm{sign}}
\xrightarrow{\mathcal P_{\mathrm{crit}}}(\phi_m,w_m)
\\ &\longrightarrow u_0\longrightarrow f_\lambda
\longrightarrow\text{retours et renormalisation}.
\end{aligned}\]

L’écriture abrégée conserve comme antécédent l’état complet et ses prolongements ; la famille scalaire est une lecture de cet antécédent. Le lecteur déjà fixé, avec son ordre, son orientation et son secteur uniforme, est pris comme structure de départ explicite. Les conclusions portent sur ce lecteur précis ; l’unicité de ce lecteur parmi toutes les lectures possibles du TPK reste hors de ces preuves.

La coordonnée \(\pi_{\mathrm{HMT}}\) utilisée par la carte appartient aux sorties internes du corpus. Son facteur s’annule exactement lors de la normalisation du deuxième moment. Les valeurs de Feigenbaum interviennent exclusivement dans la reconnaissance ultérieure, jamais dans la définition de la famille ni dans la certification des racines.


\subsection{Profil exact et criticité quadratique}
\label{hmtfg:desarrollo:2}

On conserve les notations $s_m=3m-1$, $\kappa_0=s_0$, $k_m=2s_m/\sqrt{899}$ et $f_{\lambda,0}=f_\lambda$ du profil déjà démontré. Son prolongement entier possède la série rationnelle
\[
u_0(z)=\sum_{j\ge1}(-1)^{j+1}
\frac{4^j\sum_m s_m^{2j}}{12\,899^j(2j)!}z^{2j}
=z^2-\frac{715769}{2424603}z^4
+\frac{1353832978}{32695771455}z^6+\cdots.
\]
Ainsi, le registre fini de douze secteurs produit un profil entier, dont la criticité quadratique et l’échelle sont fixées avant l’itération.


\subsection{Robustesse structurelle avec bornes uniformes}
\label{hmtfg:desarrollo:3}

La déformation est celle définie en \eqref{mab:eq:pesos-dodecafasicos-armonizados}--\eqref{mab:eq:familia-feigenbaum-ponderada-armonizada}. On écrit $M_{2,\eta}=\sum_mw_m\phi_m^2$ et $\kappa_\eta=s_\eta=\sqrt{2/M_{2,\eta}}$, en conservant les mêmes poids, phases et profils.
Les notations sont $p_m=p(\phi_m)=12\cos(3\phi_m)-\cos(4\phi_m)$ et $k_{m,\eta}=\kappa_\eta\phi_m$ ; dans les formules pondérées suivantes, $k_m$ abrège $k_{m,\eta}$, tandis que $k_{m,0}$ correspond au secteur uniforme.

\textbf{Théorème de robustesse structurelle.} Pour \(|\eta|\le1/40\), les poids sont positifs et leur somme vaut un. Le profil est pair, entier, satisfait \(u_\eta''(0)=2\) et est strictement croissant sur \((0,1]\). Pour \(0<\lambda\le2/u_\eta(1)\), l’application envoie \([-1,1]\) dans lui-même, possède un unique point critique dans cet intervalle et satisfait
\[
S f_{\lambda,\eta}(x)\le-\frac{640}{47647}<0,\qquad0<|x|\le1.
\]

\textbf{Démonstration.} Les sommes des cosinus \(3\phi_m\) et \(4\phi_m\) s’annulent par cycles de longueurs quatre et trois. Comme \(|p_m|\le13\),
\[
\frac9{160}\le w_m\le\frac{53}{480},\qquad
\frac{27}{40}M_{2,0}\le M_{2,\eta}\le\frac{53}{40}M_{2,0}.
\]
Il en découle que
\[
k_{\max,\eta}^2\le\frac{196000}{24273}<9<\pi^2,\qquad
k_{\min,\eta}^2\ge\frac{640}{47647}.
\]
Pour \(0<x\le1\), tous les sinus \(\sin(k_{m,\eta}x)\) sont positifs, d’où
\[
u_\eta'(x)=\sum_mw_mk_m\sin(k_mx)>0,\qquad
u_\eta'''(x)=-\sum_mw_mk_m^3\sin(k_mx)<0.
\]
Le quotient \(u_\eta'''/u_\eta'\) est l’opposé d’une moyenne positive des \(k_m^2\). L’invariance de la schwarzienne sous les transformations affines de la valeur donne
\[
S f_{\lambda,\eta}
=\frac{u_\eta'''}{u_\eta'}-\frac32
\left(\frac{u_\eta''}{u_\eta'}\right)^2
\le-k_{\min,\eta}^2.
\]
La parité couvre \(x<0\), et l’image de l’intervalle est
\([1-\lambda u_\eta(1),1]\). \(\square\)

La largeur \(1/40\) est un choix permettant d’obtenir des bornes uniformes, distinct d’une constante fondamentale ou d’un seuil optimal. Le secteur original reste \(\eta=0\).


\subsection{Coordonnée holomorphe quadratique globale sur une bande}
\label{hmtfg:desarrollo:4}

\textbf{Théorème de factorisation holomorphe.} Pour \(|\eta|\le1/40\), dans
\[
\mathcal S=\{z\in\mathbb C:|\operatorname{Re}z|<\pi/3\}
\]
il existe une fonction holomorphe, impaire et univalente \(b_\eta\), avec \(b_\eta'(0)=1\), telle que
\[
\boxed{u_\eta(z)=b_\eta(z)^2,\qquad
f_{\lambda,\eta}(z)=1-\lambda b_\eta(z)^2.}
\]

\textbf{Démonstration.} Pour \(z=x+iy\) dans la demi-bande droite,
\[
\operatorname{Re}u_\eta'(z)
=\sum_mw_mk_m\sin(k_mx)\cosh(k_my)>0.
\]
L’intégrale de \(u_\eta'\) sur le segment entre deux points, divisée par la différence des extrémités, a une partie réelle positive. La convexité de la demi-bande y prouve l’injectivité. La parité donne le résultat correspondant à gauche.

De plus,
\[
\operatorname{Im}u_\eta(x+iy)
=\sum_mw_m\sin(k_mx)\sinh(k_my)
\]
a le signe de \(y\) pour \(x>0\). Sur l’axe réel positif, \(u_\eta>0\) ; sur l’axe imaginaire,
\[
u_\eta(iy)=\sum_mw_m(1-\cosh(k_my))
\]
est strictement négatif hors de zéro et décroît strictement avec \(|y|\). Ces séparations, l’injectivité de chaque demi-bande et la parité prouvent que les fibres de \(u_\eta\) dans \(\mathcal S\) sont exactement \(\{z,-z\}\), sauf la fibre de l’origine.

La fonction \(u_\eta(z)/z^2\), prolongée par la valeur un en zéro, est holomorphe et ne possède aucun zéro dans la bande simplement connexe. Elle possède un logarithme holomorphe \(L_\eta\) avec \(L_\eta(0)=0\), qui est pair par unicité de la normalisation. Définissons
\[
b_\eta(z)=z\exp\!\left(\frac12L_\eta(z)\right).
\]
Alors \(b_\eta^2=u_\eta\), \(b_\eta\) est impaire et \(b_\eta'(0)=1\). L’égalité \(b_\eta(z)=b_\eta(w)\) impose \(w=z\) ou \(w=-z\). Le second cas ne donne l’égalité qu’à l’origine. Cela prouve son univalence. \(\square\)

Les douze phases déterminent ainsi une déformation conforme explicite du pli quadratique. La conjugaison dynamique complète conserve cette déformation :
\[
b_\eta\circ f_{\lambda,\eta}\circ b_\eta^{-1}(w)
=b_\eta(1-\lambda w^2),
\]
sur leurs domaines de composition. La factorisation démontrée et une conjugaison à toute la famille quadratique sont des affirmations distinctes.


\subsection{Retour, changement d’échelle et mémoire}
\label{hmtfg:desarrollo:5}

Soient \(a=-f(1)>0\), \(H_a(x)=-ax\), \(J=H_a([-1,1])\). Lorsque \(J\subset[-1,1]\) est admissible pour le retour, \(f^2(J)\subset J\) et que les compositions sont définies,
\[
\mathcal Rf=H_a^{-1}f^2H_a,\qquad
\mathcal Rf(x)=-a^{-1}f(f(-ax)).
\]
Si, de plus, \(f(J)\subset(0,1]\), un unique maximum quadratique central est conservé :
\[
(\mathcal Rf)(0)=1,\qquad
(\mathcal Rf)''(0)=-a f'(1)f''(0)<0.
\]
La composition des schwarziennes donne, aux points réguliers,
\[
S(\mathcal Rf)(x)=a^2
\left[(Sf)(f(-ax))f'(-ax)^2+(Sf)(-ax)\right]<0.
\]
Les inclusions de domaines font partie des hypothèses et sont vérifiées à chaque échelle.

Soit \(v_\eta=(u_\eta|_{[0,1]})^{-1}\). Les branches inverses sont
\[
\beta_\sigma(y)=\sigma v_\eta((1-y)/\lambda),\quad
\sigma\in\{-1,+1\},
\]
avec le registre critique \(\beta_0(1)=0\). Le pas
\[
(x,w)\longmapsto(f(x),w\mathbin{\|}\sigma(x))
\]
s’inverse sur son image en lisant la dernière feuille et en appliquant \(\beta_\sigma\). Le mot conserve les décisions de branche ; la valeur antérieure se reconstruit par l’équation de l’application.

Le retour regroupe deux transitions. À partir de l’extrémité \(x'\) et de ses feuilles \((\sigma_0,\sigma_1)\), on récupère
\[
x=H_a^{-1}\beta_{\sigma_0}
\!\left(\beta_{\sigma_1}(H_ax')\right).
\]
Les points intermédiaires sont également reconstruits. Les registres APP–TRIT–TPK de feuille, de retenue, de route, de frontière et de mémoire sont concaténés en conservant leur identité. Le signe de la branche analytique est un registre supplémentaire : il ne remplace pas ces données par une étiquette binaire.

Si \(C_N\) regroupe les histoires de longueur \(2N\) par couples, les restrictions satisfont
\[
\pi^{\mathrm{ret}}_{N,M}C_N
=C_M\pi^f_{2N,2M},\qquad M\le N.
\]
Les deux membres suppriment les mêmes couples terminaux et conservent les mêmes branches inverses. Le transport compatible des histoires finies est ainsi démontré.


\subsection{Échelles accumulées et dérivée de l’opérateur}
\label{hmtfg:desarrollo:6}

Soient \(c_n=f^{2^n}(0)\), \(c_0=1\) et \(H_n(x)=c_nx\). Tant que les compositions sont admissibles et que \(c_n\ne0\),
\[
\boxed{\mathcal R^nf=H_n^{-1}f^{2^n}H_n},\qquad
(\mathcal R^nf)(1)=\frac{c_{n+1}}{c_n}=-a_n.
\]
La récurrence découle de \(H_nH_{a_n}=H_{n+1}\). La convention \(a_n>0\) exige l’alternance du signe de \(c_n\). À une racine superstable de période \(2^N\), \(c_N=0\) : cette normalisation devient singulière et il faut utiliser les niveaux antérieurs ou une limite contrôlée.

Pour \(z=-ax\), \(w=f(z)\) et une perturbation \(h\), la dérivée complète est
\[
\boxed{
D\mathcal R_f[h](x)=
-\frac{h(w)+f'(w)h(z)}a+
\frac{h(1)}a
\left[\mathcal Rf(x)-xf'(w)f'(z)\right].
}
\]
On dérive en utilisant \(\dot a=-h(1)\), \(\dot z=xh(1)\),
\(\dot w=h(z)+xf'(z)h(1)\) et la variation du dénominateur. Si \(h(0)=0\), on obtient \(D\mathcal R_f[h](0)=0\). La direction paramétrique initiale HMT est \(h=-u_\eta\).

Cette formule inclut la variation d’échelle ; figer \(a\) donnerait un autre opérateur linéarisé.


\subsection{Racines et itinéraires certifiés jusqu’à la période 256}
\label{hmtfg:desarrollo:7}

Soit \(F_n(\lambda)=f_{\lambda,0}^{2^n}(0)\). On calcule
\[
x_0=p_0=0,\qquad x_{j+1}=1-\lambda u_0(x_j),\qquad
p_{j+1}=-u_0(x_j)-\lambda u_0'(x_j)p_j.
\]
Alors \(F_n=x_{2^n}\) et \(F_n'=p_{2^n}\).

Le programme joint utilise des intervalles rationnels avec arrondi entier vers l’extérieur, une racine carrée encadrée par une racine entière et des sinus/cosinus avec reste de Taylor explicite. La borne exacte \(|u_0''|\le2\) permet une évaluation centrée avec reste quadratique.

Les approximations historiques proposent seulement des boîtes de rayon \(10^{-42}\). On certifie des signes opposés à leurs extrémités, \(F_n'\) séparé de zéro dans toute la boîte et l’exclusion des retours de périodes qui divisent proprement \(2^n\). Le théorème des valeurs intermédiaires et la monotonie prouvent l’existence et l’unicité locales ; l’exclusion des diviseurs prouve la période minimale.

\begingroup\small
\begin{center}\begin{tabular}{p{.12\linewidth}p{.17\linewidth}p{.62\linewidth}}
\hline
Niveau & Période minimale & Centre de la boîte, arrondi \\
\hline
1 & 2 & 1.345913990471832792210949 \\
2 & 4 & 1.763379787846910127290794 \\
3 & 8 & 1.850413796950136464530567 \\
4 & 16 & 1.868979756606798125150574 \\
5 & 32 & 1.872955034435659740004205 \\
6 & 64 & 1.873806367467030119412262 \\
7 & 128 & 1.873988695221365362307169 \\
8 & 256 & 1.874027744154450672897008 \\
\hline\end{tabular}\end{center}
\endgroup

La première racine possède l’expression exacte \(\lambda_1=1/u_0(1)\). Les boîtes complètes et leurs vérifications figurent dans le certificat JSON.

Les 502 états de préclôture de ces huit niveaux, tous séparés de zéro, ainsi que leurs mots de signes complets, ont été encadrés. En particulier,
\[
\operatorname{sign}f_{\lambda_n,0}^{2^j}(0)=(-1)^j,\qquad0\le j<n.
\]
Les orientations finies des changements d’échelle antérieurs à la clôture sont établies. Les inclusions d’intervalles restrictifs à toutes les échelles constituent une vérification supplémentaire.

Pour \(\delta_{\mathrm F,n}=(\lambda_{n-1}-\lambda_{n-2})/(\lambda_n-\lambda_{n-1})\), un intervalle décimal extérieur abrégé est
\[
\boxed{\begin{gathered}
4.66921218915020415455609621588966392804<\delta_{\mathrm F,8},\\
\delta_{\mathrm F,8}<4.66921218915020415455609621588966392863.
\end{gathered}}
\]
La largeur de l’intervalle complet est inférieure à \(5.808\cdot10^{-37}\). Elle mesure l’erreur numérique du rapport fini \(\delta_{\mathrm F,8}\), distincte de l’erreur d’approximation de la limite universelle.


\subsection{Persistance analytique des racines déformées}
\label{hmtfg:desarrollo:8}

Pour chaque niveau certifié, \(F_n'(\lambda_n)\ne0\). Le théorème des fonctions implicites produit une branche analytique locale \(\lambda_n(\eta)\) avec
\[
f_{\lambda_n(\eta),\eta}^{2^n}(0)=0.
\]
Par continuité, les retours propres restent séparés de zéro pour \(\eta\) suffisamment petit ; la période minimale et l’itinéraire sont préservés. L’existence de ces voisinages locaux est démontrée. Leur largeur commune n’est pas encore identifiée à toute la bande structurelle \(1/40\).

Avec \(M=M_{2,\eta}\), la réponse du profil est
\[
v_\eta(x)=\partial_\eta u_\eta(x)
=\sum_m\frac{p_m}{12}(1-\cos(k_{m,\eta}x))
-\frac{M'}{2M}\,x u_\eta'(x).
\]
Le second terme conserve la normalisation du deuxième moment. Si \(q_j=\partial_\eta x_j\) à \(\lambda\) fixé,
\[
q_0=0,\qquad q_{j+1}=-\lambda v_\eta(x_j)-\lambda u_\eta'(x_j)q_j,
\qquad
\boxed{\lambda_n'(\eta)=-q_{2^n}/p_{2^n}}.
\]
La chaîne poids–échelle–réponse–paramètre superstable est ainsi explicitée. Cette non-dégénérescence finie est distincte de la transversalité asymptotique à la feuille stable du point fixe universel.


\subsection{Queues analytiques et propagation de l’erreur}
\label{hmtfg:desarrollo:9}

Pour \(\eta=0\), soient
\[
\mu_{2j}=\frac1{12}\sum_mk_m^{2j},\quad
p_N(z)=\sum_{j=1}^N(-1)^{j+1}\frac{\mu_{2j}}{(2j)!}z^{2j},
\quad\Omega=\frac{70}{\sqrt{899}}.
\]
Si \(q_N=\Omega^2r^2/[(2N+3)(2N+4)]<1\),
\[
\|u_0-p_N\|_r\le
\frac{\mu_{2N+2}r^{2N+2}}{(2N+2)!(1-q_N)}.
\]
L’inégalité \(\mu_{2j+2}\le\Omega^2\mu_{2j}\) majore la queue par une série géométrique. Pour \(0\le s\le2N+2\),
\[
\|(u_0-p_N)^{(s)}\|_r\le
\frac{\mu_{2N+2}r^{2N+2-s}}{(2N+2-s)!}
\left(1-\frac{\Omega^2r^2}{(2N+3-s)(2N+4-s)}\right)^{-1},
\]
lorsque le dénominateur est positif.

Pour transporter les erreurs par un retour, on contrôle également ses domaines : si \(\|f-g\|\le\varepsilon\), \(a_f,a_g\ge a_*>0\), que les compositions restent dans un domaine convexe commun et que l’on y a \(\|f'\|\le L,\ \|g\|\le M\), alors
\[
\|\mathcal Rf-\mathcal Rg\|_r
\le\varepsilon\left[\frac{1+L+L^2r}{a_*}+\frac{M}{a_*^2}\right].
\]
La preuve additionne la perturbation extérieure, la perturbation intérieure, la modification de l’argument due à \(|a_f-a_g|\le\varepsilon\) et la variation du dénominateur. Ces bornes évitent de réutiliser la somme initiale de cosinus comme si tous les profils renormalisés conservaient cette même forme.
